% Reusable colored theorem and note environments.
% Load in the preamble with: % Reusable colored theorem and note environments.
% Load in the preamble with: % Reusable colored theorem and note environments.
% Load in the preamble with: % Reusable colored theorem and note environments.
% Load in the preamble with: \input{colored-boxes}

\RequirePackage[dvipsnames]{xcolor}
\RequirePackage{amsthm}
\RequirePackage{thmtools}
\RequirePackage[framemethod=TikZ]{mdframed}
\RequirePackage{tcolorbox}
\tcbuselibrary{breakable}

\mdfsetup{skipabove=1em,skipbelow=0.5em}
\theoremstyle{definition}

% Green box: definitions.
\declaretheoremstyle[
  headfont=\bfseries\sffamily\color{ForestGreen!70!black},
  bodyfont=\normalfont,
  mdframed={
    linewidth=2pt,
    rightline=false,
    topline=false,
    bottomline=false,
    linecolor=ForestGreen,
    backgroundcolor=ForestGreen!5,
  }
]{thmgreenbox}

% Blue box: examples and exercises.
\declaretheoremstyle[
  headfont=\bfseries\sffamily\color{NavyBlue!70!black},
  bodyfont=\normalfont,
  mdframed={
    linewidth=2pt,
    rightline=false,
    topline=false,
    bottomline=false,
    linecolor=NavyBlue,
    backgroundcolor=NavyBlue!5,
  }
]{thmbluebox}

% Blue line: remarks and notes.
\declaretheoremstyle[
  headfont=\bfseries\sffamily\color{NavyBlue!70!black},
  bodyfont=\normalfont,
  mdframed={
    linewidth=2pt,
    rightline=false,
    topline=false,
    bottomline=false,
    linecolor=NavyBlue,
  }
]{thmblueline}

% Sienna box: mathematical results.
\declaretheoremstyle[
  headfont=\bfseries\sffamily\color{RawSienna!70!black},
  bodyfont=\normalfont,
  mdframed={
    linewidth=2pt,
    rightline=false,
    topline=false,
    bottomline=false,
    linecolor=RawSienna,
    backgroundcolor=RawSienna!5,
  }
]{thmredbox}

% Pale sienna box: proofs and solutions.
\declaretheoremstyle[
  headfont=\bfseries\sffamily\color{RawSienna!70!black},
  bodyfont=\normalfont,
  numbered=no,
  mdframed={
    linewidth=2pt,
    rightline=false,
    topline=false,
    bottomline=false,
    linecolor=RawSienna,
    backgroundcolor=RawSienna!1,
  },
  qed=\qedsymbol,
]{thmproofbox}

% Pale blue box: intuition and explanatory arguments.
\declaretheoremstyle[
  headfont=\bfseries\sffamily\color{NavyBlue!70!black},
  bodyfont=\normalfont,
  numbered=no,
  mdframed={
    linewidth=2pt,
    rightline=false,
    topline=false,
    bottomline=false,
    linecolor=NavyBlue,
    backgroundcolor=NavyBlue!1,
  }
]{thmexplanationbox}

% All numbered environments share a counter that resets each lecture section.
\declaretheorem[style=thmredbox,name=Theorem,numberwithin=section]{theorem}
\declaretheorem[style=thmgreenbox,name=Definition,sibling=theorem]{definition}
\declaretheorem[style=thmredbox,name=Proposition,sibling=theorem]{proposition}
\declaretheorem[style=thmredbox,name=Proposition,sibling=theorem]{prop}
\declaretheorem[style=thmredbox,name=Lemma,sibling=theorem]{lemma}
\declaretheorem[style=thmredbox,name=Corollary,sibling=theorem]{corollary}
\declaretheorem[style=thmbluebox,name=Exercise,sibling=theorem]{exercise}

% Unnumbered supporting environments.
\declaretheorem[style=thmbluebox,numbered=no,name=Example]{example}
\declaretheorem[style=thmbluebox,numbered=no,name=Example]{eg}
\declaretheorem[style=thmblueline,numbered=no,name=Remark]{remark}
\declaretheorem[style=thmblueline,numbered=no,name=Note]{note}

% Replace the standard proof with a colored proof box.
\declaretheorem[style=thmproofbox,name=Proof]{replacementproof}
\renewenvironment{proof}[1][\proofname]
  {\vspace{-10pt}\begin{replacementproof}}
  {\end{replacementproof}}

\declaretheorem[style=thmproofbox,name=Solution]{solution}
\declaretheorem[style=thmexplanationbox,name=Explanation]{tmpexplanation}
\newenvironment{explanation}[1][]
  {\vspace{-10pt}\begin{tmpexplanation}}
  {\end{tmpexplanation}}

% General-purpose tcolorbox notes.
\newenvironment{verbetering}{%
  \begin{tcolorbox}[
    arc=0mm,
    colback=white,
    colframe=green!60!black,
    title=Opmerking,
    fonttitle=\sffamily,
    breakable,
  ]
}{\end{tcolorbox}}

\newenvironment{noot}[1]{%
  \begin{tcolorbox}[
    arc=0mm,
    colback=white,
    colframe=black!40,
    title=#1,
    fonttitle=\sffamily,
    breakable,
  ]
}{\end{tcolorbox}}

\RequirePackage[dvipsnames]{xcolor}
\RequirePackage{amsthm}
\RequirePackage{thmtools}
\RequirePackage[framemethod=TikZ]{mdframed}
\RequirePackage{tcolorbox}
\tcbuselibrary{breakable}

\mdfsetup{skipabove=1em,skipbelow=0.5em}
\theoremstyle{definition}

% Green box: definitions.
\declaretheoremstyle[
  headfont=\bfseries\sffamily\color{ForestGreen!70!black},
  bodyfont=\normalfont,
  mdframed={
    linewidth=2pt,
    rightline=false,
    topline=false,
    bottomline=false,
    linecolor=ForestGreen,
    backgroundcolor=ForestGreen!5,
  }
]{thmgreenbox}

% Blue box: examples and exercises.
\declaretheoremstyle[
  headfont=\bfseries\sffamily\color{NavyBlue!70!black},
  bodyfont=\normalfont,
  mdframed={
    linewidth=2pt,
    rightline=false,
    topline=false,
    bottomline=false,
    linecolor=NavyBlue,
    backgroundcolor=NavyBlue!5,
  }
]{thmbluebox}

% Blue line: remarks and notes.
\declaretheoremstyle[
  headfont=\bfseries\sffamily\color{NavyBlue!70!black},
  bodyfont=\normalfont,
  mdframed={
    linewidth=2pt,
    rightline=false,
    topline=false,
    bottomline=false,
    linecolor=NavyBlue,
  }
]{thmblueline}

% Sienna box: mathematical results.
\declaretheoremstyle[
  headfont=\bfseries\sffamily\color{RawSienna!70!black},
  bodyfont=\normalfont,
  mdframed={
    linewidth=2pt,
    rightline=false,
    topline=false,
    bottomline=false,
    linecolor=RawSienna,
    backgroundcolor=RawSienna!5,
  }
]{thmredbox}

% Pale sienna box: proofs and solutions.
\declaretheoremstyle[
  headfont=\bfseries\sffamily\color{RawSienna!70!black},
  bodyfont=\normalfont,
  numbered=no,
  mdframed={
    linewidth=2pt,
    rightline=false,
    topline=false,
    bottomline=false,
    linecolor=RawSienna,
    backgroundcolor=RawSienna!1,
  },
  qed=\qedsymbol,
]{thmproofbox}

% Pale blue box: intuition and explanatory arguments.
\declaretheoremstyle[
  headfont=\bfseries\sffamily\color{NavyBlue!70!black},
  bodyfont=\normalfont,
  numbered=no,
  mdframed={
    linewidth=2pt,
    rightline=false,
    topline=false,
    bottomline=false,
    linecolor=NavyBlue,
    backgroundcolor=NavyBlue!1,
  }
]{thmexplanationbox}

% All numbered environments share a counter that resets each lecture section.
\declaretheorem[style=thmredbox,name=Theorem,numberwithin=section]{theorem}
\declaretheorem[style=thmgreenbox,name=Definition,sibling=theorem]{definition}
\declaretheorem[style=thmredbox,name=Proposition,sibling=theorem]{proposition}
\declaretheorem[style=thmredbox,name=Proposition,sibling=theorem]{prop}
\declaretheorem[style=thmredbox,name=Lemma,sibling=theorem]{lemma}
\declaretheorem[style=thmredbox,name=Corollary,sibling=theorem]{corollary}
\declaretheorem[style=thmbluebox,name=Exercise,sibling=theorem]{exercise}

% Unnumbered supporting environments.
\declaretheorem[style=thmbluebox,numbered=no,name=Example]{example}
\declaretheorem[style=thmbluebox,numbered=no,name=Example]{eg}
\declaretheorem[style=thmblueline,numbered=no,name=Remark]{remark}
\declaretheorem[style=thmblueline,numbered=no,name=Note]{note}

% Replace the standard proof with a colored proof box.
\declaretheorem[style=thmproofbox,name=Proof]{replacementproof}
\renewenvironment{proof}[1][\proofname]
  {\vspace{-10pt}\begin{replacementproof}}
  {\end{replacementproof}}

\declaretheorem[style=thmproofbox,name=Solution]{solution}
\declaretheorem[style=thmexplanationbox,name=Explanation]{tmpexplanation}
\newenvironment{explanation}[1][]
  {\vspace{-10pt}\begin{tmpexplanation}}
  {\end{tmpexplanation}}

% General-purpose tcolorbox notes.
\newenvironment{verbetering}{%
  \begin{tcolorbox}[
    arc=0mm,
    colback=white,
    colframe=green!60!black,
    title=Opmerking,
    fonttitle=\sffamily,
    breakable,
  ]
}{\end{tcolorbox}}

\newenvironment{noot}[1]{%
  \begin{tcolorbox}[
    arc=0mm,
    colback=white,
    colframe=black!40,
    title=#1,
    fonttitle=\sffamily,
    breakable,
  ]
}{\end{tcolorbox}}

\RequirePackage[dvipsnames]{xcolor}
\RequirePackage{amsthm}
\RequirePackage{thmtools}
\RequirePackage[framemethod=TikZ]{mdframed}
\RequirePackage{tcolorbox}
\tcbuselibrary{breakable}

\mdfsetup{skipabove=1em,skipbelow=0.5em}
\theoremstyle{definition}

% Green box: definitions.
\declaretheoremstyle[
  headfont=\bfseries\sffamily\color{ForestGreen!70!black},
  bodyfont=\normalfont,
  mdframed={
    linewidth=2pt,
    rightline=false,
    topline=false,
    bottomline=false,
    linecolor=ForestGreen,
    backgroundcolor=ForestGreen!5,
  }
]{thmgreenbox}

% Blue box: examples and exercises.
\declaretheoremstyle[
  headfont=\bfseries\sffamily\color{NavyBlue!70!black},
  bodyfont=\normalfont,
  mdframed={
    linewidth=2pt,
    rightline=false,
    topline=false,
    bottomline=false,
    linecolor=NavyBlue,
    backgroundcolor=NavyBlue!5,
  }
]{thmbluebox}

% Blue line: remarks and notes.
\declaretheoremstyle[
  headfont=\bfseries\sffamily\color{NavyBlue!70!black},
  bodyfont=\normalfont,
  mdframed={
    linewidth=2pt,
    rightline=false,
    topline=false,
    bottomline=false,
    linecolor=NavyBlue,
  }
]{thmblueline}

% Sienna box: mathematical results.
\declaretheoremstyle[
  headfont=\bfseries\sffamily\color{RawSienna!70!black},
  bodyfont=\normalfont,
  mdframed={
    linewidth=2pt,
    rightline=false,
    topline=false,
    bottomline=false,
    linecolor=RawSienna,
    backgroundcolor=RawSienna!5,
  }
]{thmredbox}

% Pale sienna box: proofs and solutions.
\declaretheoremstyle[
  headfont=\bfseries\sffamily\color{RawSienna!70!black},
  bodyfont=\normalfont,
  numbered=no,
  mdframed={
    linewidth=2pt,
    rightline=false,
    topline=false,
    bottomline=false,
    linecolor=RawSienna,
    backgroundcolor=RawSienna!1,
  },
  qed=\qedsymbol,
]{thmproofbox}

% Pale blue box: intuition and explanatory arguments.
\declaretheoremstyle[
  headfont=\bfseries\sffamily\color{NavyBlue!70!black},
  bodyfont=\normalfont,
  numbered=no,
  mdframed={
    linewidth=2pt,
    rightline=false,
    topline=false,
    bottomline=false,
    linecolor=NavyBlue,
    backgroundcolor=NavyBlue!1,
  }
]{thmexplanationbox}

% All numbered environments share a counter that resets each lecture section.
\declaretheorem[style=thmredbox,name=Theorem,numberwithin=section]{theorem}
\declaretheorem[style=thmgreenbox,name=Definition,sibling=theorem]{definition}
\declaretheorem[style=thmredbox,name=Proposition,sibling=theorem]{proposition}
\declaretheorem[style=thmredbox,name=Proposition,sibling=theorem]{prop}
\declaretheorem[style=thmredbox,name=Lemma,sibling=theorem]{lemma}
\declaretheorem[style=thmredbox,name=Corollary,sibling=theorem]{corollary}
\declaretheorem[style=thmbluebox,name=Exercise,sibling=theorem]{exercise}

% Unnumbered supporting environments.
\declaretheorem[style=thmbluebox,numbered=no,name=Example]{example}
\declaretheorem[style=thmbluebox,numbered=no,name=Example]{eg}
\declaretheorem[style=thmblueline,numbered=no,name=Remark]{remark}
\declaretheorem[style=thmblueline,numbered=no,name=Note]{note}

% Replace the standard proof with a colored proof box.
\declaretheorem[style=thmproofbox,name=Proof]{replacementproof}
\renewenvironment{proof}[1][\proofname]
  {\vspace{-10pt}\begin{replacementproof}}
  {\end{replacementproof}}

\declaretheorem[style=thmproofbox,name=Solution]{solution}
\declaretheorem[style=thmexplanationbox,name=Explanation]{tmpexplanation}
\newenvironment{explanation}[1][]
  {\vspace{-10pt}\begin{tmpexplanation}}
  {\end{tmpexplanation}}

% General-purpose tcolorbox notes.
\newenvironment{verbetering}{%
  \begin{tcolorbox}[
    arc=0mm,
    colback=white,
    colframe=green!60!black,
    title=Opmerking,
    fonttitle=\sffamily,
    breakable,
  ]
}{\end{tcolorbox}}

\newenvironment{noot}[1]{%
  \begin{tcolorbox}[
    arc=0mm,
    colback=white,
    colframe=black!40,
    title=#1,
    fonttitle=\sffamily,
    breakable,
  ]
}{\end{tcolorbox}}

\RequirePackage[dvipsnames]{xcolor}
\RequirePackage{amsthm}
\RequirePackage{thmtools}
\RequirePackage[framemethod=TikZ]{mdframed}
\RequirePackage{tcolorbox}
\tcbuselibrary{breakable}

\mdfsetup{skipabove=1em,skipbelow=0.5em}
\theoremstyle{definition}

% Green box: definitions.
\declaretheoremstyle[
  headfont=\bfseries\sffamily\color{ForestGreen!70!black},
  bodyfont=\normalfont,
  mdframed={
    linewidth=2pt,
    rightline=false,
    topline=false,
    bottomline=false,
    linecolor=ForestGreen,
    backgroundcolor=ForestGreen!5,
  }
]{thmgreenbox}

% Blue box: examples and exercises.
\declaretheoremstyle[
  headfont=\bfseries\sffamily\color{NavyBlue!70!black},
  bodyfont=\normalfont,
  mdframed={
    linewidth=2pt,
    rightline=false,
    topline=false,
    bottomline=false,
    linecolor=NavyBlue,
    backgroundcolor=NavyBlue!5,
  }
]{thmbluebox}

% Blue line: remarks and notes.
\declaretheoremstyle[
  headfont=\bfseries\sffamily\color{NavyBlue!70!black},
  bodyfont=\normalfont,
  mdframed={
    linewidth=2pt,
    rightline=false,
    topline=false,
    bottomline=false,
    linecolor=NavyBlue,
  }
]{thmblueline}

% Sienna box: mathematical results.
\declaretheoremstyle[
  headfont=\bfseries\sffamily\color{RawSienna!70!black},
  bodyfont=\normalfont,
  mdframed={
    linewidth=2pt,
    rightline=false,
    topline=false,
    bottomline=false,
    linecolor=RawSienna,
    backgroundcolor=RawSienna!5,
  }
]{thmredbox}

% Pale sienna box: proofs and solutions.
\declaretheoremstyle[
  headfont=\bfseries\sffamily\color{RawSienna!70!black},
  bodyfont=\normalfont,
  numbered=no,
  mdframed={
    linewidth=2pt,
    rightline=false,
    topline=false,
    bottomline=false,
    linecolor=RawSienna,
    backgroundcolor=RawSienna!1,
  },
  qed=\qedsymbol,
]{thmproofbox}

% Pale blue box: intuition and explanatory arguments.
\declaretheoremstyle[
  headfont=\bfseries\sffamily\color{NavyBlue!70!black},
  bodyfont=\normalfont,
  numbered=no,
  mdframed={
    linewidth=2pt,
    rightline=false,
    topline=false,
    bottomline=false,
    linecolor=NavyBlue,
    backgroundcolor=NavyBlue!1,
  }
]{thmexplanationbox}

% All numbered environments share a counter that resets each lecture section.
\declaretheorem[style=thmredbox,name=Theorem,numberwithin=section]{theorem}
\declaretheorem[style=thmgreenbox,name=Definition,sibling=theorem]{definition}
\declaretheorem[style=thmredbox,name=Proposition,sibling=theorem]{proposition}
\declaretheorem[style=thmredbox,name=Proposition,sibling=theorem]{prop}
\declaretheorem[style=thmredbox,name=Lemma,sibling=theorem]{lemma}
\declaretheorem[style=thmredbox,name=Corollary,sibling=theorem]{corollary}
\declaretheorem[style=thmbluebox,name=Exercise,sibling=theorem]{exercise}

% Unnumbered supporting environments.
\declaretheorem[style=thmbluebox,numbered=no,name=Example]{example}
\declaretheorem[style=thmbluebox,numbered=no,name=Example]{eg}
\declaretheorem[style=thmblueline,numbered=no,name=Remark]{remark}
\declaretheorem[style=thmblueline,numbered=no,name=Note]{note}

% Replace the standard proof with a colored proof box.
\declaretheorem[style=thmproofbox,name=Proof]{replacementproof}
\renewenvironment{proof}[1][\proofname]
  {\vspace{-10pt}\begin{replacementproof}}
  {\end{replacementproof}}

\declaretheorem[style=thmproofbox,name=Solution]{solution}
\declaretheorem[style=thmexplanationbox,name=Explanation]{tmpexplanation}
\newenvironment{explanation}[1][]
  {\vspace{-10pt}\begin{tmpexplanation}}
  {\end{tmpexplanation}}

% General-purpose tcolorbox notes.
\newenvironment{verbetering}{%
  \begin{tcolorbox}[
    arc=0mm,
    colback=white,
    colframe=green!60!black,
    title=Opmerking,
    fonttitle=\sffamily,
    breakable,
  ]
}{\end{tcolorbox}}

\newenvironment{noot}[1]{%
  \begin{tcolorbox}[
    arc=0mm,
    colback=white,
    colframe=black!40,
    title=#1,
    fonttitle=\sffamily,
    breakable,
  ]
}{\end{tcolorbox}}